El núcleo adaptativo del sistema híbrido reside en la capacidad de inferir, a partir de variables operativas medidas en tiempo real, la biomasa larval activa que alimenta el modelo DEB. Dado que la medición destructiva de $B(t)$ es esporádica e invasiva, se recurre a un estimador externo no paramétrico basado en Procesos Gaussianos (GP). A diferencia de un observador de estado clásico (e.g., filtro de Kalman) que exige un modelo de transición explícito para la variable oculta, el GP opera como un regresor probabilístico que aprende directamente el mapeo $\mathbf{x}(t) \mapsto \hat{B}(t)$ desde datos empíricos, sin imponer una dinámica a priori sobre $B(t)$. Esta elección preserva la integridad del modelo DEB —que ya contiene la dinámica bioenergética— y evita la circularidad de estimar $B(t)$ usando las mismas EDO que luego la consumen \parencite{raissiPhysicsinformedNeuralNetworks2019, weichertReviewMachineLearning2019}.

\subsection{Justificación fuera de la ODE y criterio de parsimonia}
\label{subsec:cap4_gp_justificacion}

La ubicación del estimador \emph{fuera} del sistema de ecuaciones diferenciales responde a una decisión epistémica deliberada. Si $B(t)$ se tratara como un estado oculto dentro de un observador tipo Kalman-Bucy, sería necesario postular una ecuación de transición adicional para la biomasa —por ejemplo, una extensión del propio DEB— lo que redundaría con la estructura mecanicista ya establecida en el \cref{chap:cap3}. Más aún, un filtro de Kalman linealizado (EKF) requeriría la linealización del sistema DEB alrededor de la trayectoria nominal, introduciendo errores de aproximación en un régimen altamente no lineal (cinética de Holling II, modulación térmica tipo Sharpe-Schoolfield).

El GP, en cambio, se sitúa en una capa independiente que consume los \emph{residuos de información} contenidos en las variables operativas: la concentración de \ce{CO2} acumulada refleja la historia respiratoria del reactor, la temperatura codifica la actividad metabólica exotérmica, y la humedad relativa está correlacionada con la degradación del sustrato y la densidad larval efectiva \parencite{rossiEstimatingDynamicsGreenhouse2024}. Al desacoplar la inferencia de la integración numérica, el sistema mantiene la interpretabilidad física del DEB (cada término de las EDO conserva su significado bioquímico) mientras incorpora flexibilidad empírica para compensar perturbaciones no modeladas.

El criterio de parsimonia impulsa la elección del GP frente a arquitecturas de aprendizaje profundo más complejas. En el contexto experimental de esta tesis, un ciclo completo de bioconversión genera entre \num{3} y \num{5} mediciones destructivas de biomasa por réplica, con un número total de réplicas típicamente inferior a \num{6}. Un modelo de redes neuronales profundas, con miles de parámetros libres, inevitablemente sobreajustaría un conjunto de entrenamiento de $N \approx \num{30}$ observaciones. En contraste, un GP con kernel de covarianza de base radial (RBF) y escalas de longitud automáticas (ARD) se parametriza mediante un vector de hiperparámetros de dimensión $D+2$ (donde $D=4$ es el número de entradas), resultando en un modelo altamente regularizado por la estructura del prior bayesiano \parencite{adadiPeekingBlackBoxSurvey2018, biagiDevelopmentMachineLearningbased2024}.

\subsection{Variables de entrada y codificación del estado del reactor}
\label{subsec:cap4_gp_variables}

El vector de entrada al estimador GP se define como
\begin{equation}
    \label{eq:cap4_vector_entrada}
    \mathbf{x}(t) = \bigl[\, C(t),\; T(t),\; \mathrm{HR}(t),\; t \,\bigr]^\intercal \in \mathbb{R}^4,
\end{equation}
donde cada componente codifica un aspecto distinto del estado fisiológico del reactor:

\begin{itemize}
    \item $C(t)$: concentración de \ce{CO2} en ppm, medida por sensor NDIR. Actúa como proxy integrado de la respiración acumulada. Dado que la tasa de producción de \ce{CO2} es proporcional a la biomasa activa bajo condiciones aeróbicas (véase \cref{sec:cap3_formulacion_ode}), la propia concentración ambiental —en sistemas cerrados o semicerrados— guarda correlación monotónica con la masa larval presente, especialmente tras compensar por el volumen de cabeza del reactor.
    
    \item $T(t)$: temperatura del sustrato en °C. Las larvas de \textit{H. illucens} son organismos exotérmicos cuya actividad metabólica genera calor de fermentación y respiración. La temperatura del lecho refleja, con cierta inercia térmica, la densidad metabólica instantánea y, por ende, la biomasa activa \parencite{bekkerImpactSubstrateMoisture2021}.
    
    \item $\mathrm{HR}(t)$: humedad relativa en \%. La humedad intersticial del sustrato está directamente modulada por la actividad larval (fragmentación mecánica, exudación) y por la degradación microbiana. Rangos extremos de HR ($> \num{85}$\% o $< \num{40}$\%) inhiben el crecimiento, por lo que esta variable aporta información sobre la viabilidad del sistema.
    
    \item $t$: tiempo transcurrido desde el inicio del ciclo en días. Aunque el DEB ya incorpora la dinámica temporal, la inclusión explícita de $t$ en $\mathbf{x}(t)$ permite al GP capturar tendencias de crecimiento no explicadas por las variables físicas instantáneas —por ejemplo, efectos de maduración larval o cambios en la eficiencia de asimilación a lo largo del ciclo— sin modificar la estructura del modelo mecanicista.
\end{itemize}

Previo al entrenamiento, cada componente de $\mathbf{x}$ se estandariza mediante $z$-score:
\begin{equation}
    \label{eq:cap4_estandarizacion}
    \tilde{x}_d = \frac{x_d - \mu_{x_d}}{\sigma_{x_d}}, \qquad d = 1,\dots,4,
\end{equation}
donde $\mu_{x_d}$ y $\sigma_{x_d}$ se calculan sobre el conjunto de entrenamiento. Esta normalización es esencial porque las escalas naturales difieren en órdenes de magnitud (ppm vs. °C vs. días), y el kernel RBF-ARD es sensible a la métrica euclidiana en el espacio de entradas.

\subsection{Formulación del Proceso Gaussiano}
\label{subsec:cap4_gp_formulacion}

Sea $\mathcal{D} = \{(\mathbf{x}_i, B_i)\}_{i=1}^{N}$ el conjunto de entrenamiento, donde $B_i$ representa la biomasa húmeda total del reactor (en g) medida destructivamente en el instante $t_i$, y $\mathbf{x}_i = \mathbf{x}(t_i)$ el vector de entradas correspondiente. Se asume un prior gaussiano sobre la función latente $f(\mathbf{x})$:
\begin{equation}
    \label{eq:cap4_gp_prior}
    f(\mathbf{x}) \sim \mathcal{GP}\bigl(m(\mathbf{x}),\, k(\mathbf{x}, \mathbf{x}')\bigr),
\end{equation}
con función media $m(\mathbf{x}) \equiv 0$ (tras centrar las observaciones de $B$) y función de covarianza (kernel) de base radial con escalas de longitud automáticas por dimensión (ARD):
\begin{equation}
    \label{eq:cap4_kernel_rbf_ard}
    k(\mathbf{x}, \mathbf{x}') = \sigma_f^{2}\,
    \exp\!\left(-\frac{1}{2}\sum_{d=1}^{4}\frac{(x_d - x'_d)^{2}}{\ell_d^{2}}\right),
\end{equation}
donde $\sigma_f^{2}$ es la varianza de señal y $\ell_d > 0$ la escala de longitud característica de la $d$-ésima entrada. La inclusión de ARD permite que el GP aprenda, de los datos, qué variables operativas son más informativas para la predicción de biomasa: una $\ell_d$ pequeña indica alta sensibilidad a la variable $d$, mientras que $\ell_d \to \infty$ la desactiva efectivamente \parencite{raissiPhysicsinformedNeuralNetworks2019}.

El modelo de observación incorpora ruido gaussiano aditivo:
\begin{equation}
    \label{eq:cap4_modelo_ruido}
    B_i = f(\mathbf{x}_i) + \varepsilon_i, \qquad \varepsilon_i \sim \mathcal{N}(0, \sigma_n^{2}),
\end{equation}
con varianza de ruido $\sigma_n^{2}$. La matriz de covarianza del conjunto de entrenamiento resulta
\begin{equation}
    \label{eq:cap4_matriz_covarianza}
    \mathbf{K} = \mathbf{K}_{ff} + \sigma_n^{2}\mathbf{I}_N \in \mathbb{R}^{N \times N},
    \quad \text{con } [\mathbf{K}_{ff}]_{ij} = k(\mathbf{x}_i, \mathbf{x}_j).
\end{equation}

Para un punto de consulta $\mathbf{x}_*$, la distribución predictiva posterior del GP es gaussiana:
\begin{align}
    \label{eq:cap4_media_posterior}
    \mu(\mathbf{x}_*) &= \mathbf{k}_{*}^{\top}\,\mathbf{K}^{-1}\,\mathbf{y}, \\
    \label{eq:cap4_varianza_posterior}
    \sigma^{2}(\mathbf{x}_*) &= k(\mathbf{x}_*, \mathbf{x}_*) - \mathbf{k}_{*}^{\top}\,\mathbf{K}^{-1}\,\mathbf{k}_{*},
\end{align}
donde $\mathbf{y} = [B_1, \dots, B_N]^\intercal$, $\mathbf{k}_{*} = [k(\mathbf{x}_*, \mathbf{x}_1), \dots, k(\mathbf{x}_*, \mathbf{x}_N)]^\intercal$, y $k(\mathbf{x}_*, \mathbf{x}_*) = \sigma_f^{2}$. El valor esperado $\mu(\mathbf{x}_*)$ constituye el pseudo-observable $\hat{B}(t)$ que alimenta el modelo DEB, mientras que $\sigma^{2}(\mathbf{x}_*)$ cuantifica la incertidumbre epistémica asociada a la estimación.

Los hiperparámetros $\boldsymbol{\theta} = (\sigma_f^{2}, \sigma_n^{2}, \ell_1, \ell_2, \ell_3, \ell_4)$ se estiman maximizando la log-verosimilitud marginal:
\begin{equation}
    \label{eq:cap4_log_verosimilitud}
    \log p(\mathbf{y} \mid \mathbf{X}, \boldsymbol{\theta}) =
    -\frac{1}{2}\mathbf{y}^{\top}\mathbf{K}^{-1}\mathbf{y}
    -\frac{1}{2}\log|\mathbf{K}|
    -\frac{N}{2}\log(2\pi),
\end{equation}
mediante un optimizador de gradiente de segundo orden (L-BFGS-B). La maximización de esta función balancea el ajuste a los datos (primer término) contra la complejidad del modelo (segundo término), implementando de facto un principio de parsimonia bayesiana similar al criterio de información de Akaike \parencite{adadiPeekingBlackBoxSurvey2018}.

\subsection{Entrenamiento supervisado y validación leave-one-replica-out}
\label{subsec:cap4_gp_entrenamiento}

La estrategia de validación debe respetar la estructura correlacional de los datos experimentales. Las mediciones destructivas dentro de una misma réplica (mismo reactor, mismo lote de larvas, mismo sustrato) están serialmente correlacionadas; por tanto, una validación cruzada aleatoria (\emph{k-fold} estándar) violaría el supuesto de independencia y produciría estimaciones optimistas del error. En su lugar, se adopta un esquema \emph{leave-one-replica-out} (LORO): en cada iteración, se retira una réplica completa del conjunto de entrenamiento, se entrena el GP con las réplicas restantes, y se evalúa la predicción sobre todos los puntos de la réplica excluida.

El procedimiento formal es el siguiente:

\begin{enumerate}
    \item Sea $R$ el número total de réplicas experimentales ($R \geq 3$). Para cada réplica $r \in \{1, \dots, R\}$:
    \begin{enumerate}
        \item Definir $\mathcal{D}_{\setminus r} = \mathcal{D} \setminus \{(\mathbf{x}_i, B_i) : i \in \text{réplica } r\}$.
        \item Entrenar hiperparámetros $\hat{\boldsymbol{\theta}}_{-r}$ maximizando \cref{eq:cap4_log_verosimilitud} sobre $\mathcal{D}_{\setminus r}$.
        \item Predecir $\{\mu_{-r}(\mathbf{x}_i), \sigma_{-r}^{2}(\mathbf{x}_i)\}$ para todo $\mathbf{x}_i$ de la réplica $r$.
    \end{enumerate}
    \item Agregar las predicciones y calcular métricas globales de desempeño.
\end{enumerate}

Las métricas de validación seleccionadas son:
\begin{align}
    \label{eq:cap4_rmse}
    \mathrm{RMSE} &= \sqrt{\frac{1}{N}\sum_{i=1}^{N}\bigl(B_i - \mu(\mathbf{x}_i)\bigr)^{2}}, \\
    \label{eq:cap4_mape}
    \mathrm{MAPE} &= \frac{100}{N}\sum_{i=1}^{N}
    \left|\frac{B_i - \mu(\mathbf{x}_i)}{B_i}\right|, \\
    \label{eq:cap4_r2}
    R^{2} &= 1 - \frac{\sum_{i}(B_i - \mu(\mathbf{x}_i))^{2}}
    {\sum_{i}(B_i - \bar{B})^{2}}.
\end{align}

Un GP bien calibrado debe exhibir $R^{2} \geq \num{0.70}$ y un MAPE inferior al \num{15}\%, umbrales coherentes con la precisión instrumental del sensor NDIR establecida en la hipótesis H1 (véase \cref{chap:marco_metodologico}). Adicionalmente, se verifica la calibración probabilística mediante el \emph{negative log predictive density} (NLPD):
\begin{equation}
    \label{eq:cap4_nlpd}
    \mathrm{NLPD} = -\frac{1}{N}\sum_{i=1}^{N}
    \log p\bigl(B_i \mid \mu(\mathbf{x}_i), \sigma^{2}(\mathbf{x}_i)\bigr),
\end{equation}
donde $p(\cdot)$ es la densidad gaussiana evaluada en el valor observado. Un NLPD bajo indica que la varianza predictiva $\sigma^{2}(\mathbf{x}_i)$ es consistente con el error residual real, confirmando que el GP no es ni sobreconfiado ni excesivamente conservador \parencite{biagiDevelopmentMachineLearningbased2024}.

\subsection{Propagación de la incertidumbre al pseudo-observable $\hat{B}$}
\label{subsec:cap4_gp_propagacion}

La salida del estimador GP no es un valor puntual, sino una distribución completa $p(\hat{B}(t)) = \mathcal{N}(\mu(\mathbf{x}(t)), \sigma^{2}(\mathbf{x}(t)))$. Esta incertidumbre debe propagarse al modelo DEB para que la predicción de emisiones refleje fielmente la indeterminación en la variable de estado. Se distinguen dos niveles de propagación.

\paragraph{Propagación de primer orden.}
El modelo DEB predice la tasa de emisión $\dot{m}_{\mathrm{CO_2}}^{\mathrm{DEB}}(t)$ como función implícita del estado, $\dot{m}(t) = g(B(t), S(t), L(t); \boldsymbol{\theta}_{\mathrm{DEB}})$. Para una trayectoria dada, la sensibilidad de la emisión respecto a la biomasa se aproxima linealizando alrededor de $\hat{B}(t)$:
\begin{equation}
    \label{eq:cap4_propagacion_lineal}
    \sigma^{2}_{\dot{m}}(t) \approx
    \left(
    \frac{\partial \dot{m}_{\mathrm{CO_2}}^{\mathrm{DEB}}}{\partial B}
    \Big|_{B=\hat{B}(t)}
    \right)^{\!2}
    \sigma^{2}_{\hat{B}}(t).
\end{equation}
La derivada parcial se evalúa numéricamente mediante diferencias finitas durante la integración de las EDO, o analíticamente si se deriva la expresión de la respiración de mantenimiento $\dot{m}_{\mathrm{CO_2}}^{\mathrm{maint}} = \kappa\,m\,B(t)$, donde la dependencia es directamente proporcional. Esta aproximación de primer orden es válida cuando $\sigma_{\hat{B}}(t)$ es pequeña respecto a $\hat{B}(t)$, condición que se verifica a posteriori en las ventanas de validación.

\paragraph{Propagación mediante simulación.}
Cuando la varianza del GP es elevada —típicamente en los extremos del dominio temporal (inicio y final del ciclo) donde los datos de entrenamiento son escasos— la linealización puede subestimar la incertidumbre. En tales casos, se implementa un esquema Monte Carlo: se muestrean $M = \num{500}$ realizaciones $B^{(j)}(t) \sim \mathcal{N}(\hat{B}(t), \sigma^{2}_{\hat{B}}(t))$, se integra el sistema DEB para cada realización manteniendo las demás condiciones iniciales fijas, y se agrega la distribución empírica de $\dot{m}_{\mathrm{CO_2}}^{\mathrm{DEB},(j)}(t)$. El intervalo de predicción del \num{95}\% se obtiene directamente de los percentiles \num{2.5} y \num{97.5} de la muestra.

\paragraph{Definición formal del pseudo-observable.}
El estimador GP transforma las variables operativas medidas en una variable de estado no observable directamente por sensores. Se define formalmente el pseudo-observable como
\begin{equation}
    \label{eq:cap4_pseudo_observable}
    \hat{B}(t) \equiv \mathbb{E}_{\mathrm{GP}}\bigl[B \mid \mathbf{x}(t), \mathcal{D}\bigr] = \mu(\mathbf{x}(t)),
\end{equation}
con varianza asociada $\sigma^{2}_{\hat{B}}(t) = \sigma^{2}(\mathbf{x}(t))$. Al inyectar $\hat{B}(t)$ en las EDO del \cref{chap:cap3}, el sistema DEB opera en modo de \emph{lazo cerrado empírico}: la trayectoria de emisiones se ajusta dinámicamente a la realidad biológica inferida, en lugar de depender de una curva de crecimiento teórica fijada a priori. Esta retroalimentación, aunque diferida en el tiempo, es la que permite al sistema híbrido alcanzar el objetivo de MAPE $< \num{10}$\% en la predicción de emisiones, conforme a la hipótesis H3.
